Many regulatory documents specify requirements that valid data must satisfy, but those specifications are rarely available as data quality constraints. 
%%% texto 'antigo'
%When constraints are missing, the standard solution is constraint discovery, which discovers patterns in the data but cannot recover requirements derived from governing text. 
%When constraints are undefined, a well-studied solution is constraint discovery, but discovery finds patterns the data happens to satisfy, not the requirements it is supposed to satisfy. 
Data profiling, particularly constraint discovery, is a well-studied alternative when constraints are unavailable, but discovered constraints reflect what the data \add{(approximately) }satisfies rather than what it is \subs{supposed}{required} to satisfy.
We introduce the \emph{Text-to-Constraints} problem: given a dataset and related regulatory text, \subs{extract executable data quality constraints grounded in that text}{map the requirements stated in that text to executable data quality constraints over the dataset schema}. 
Then, we present the first framework to address this problem, \system, which identifies verifiable requirements in regulatory documents and maps them to denial constraints using a rich schema-grounded predicate vocabulary. 
We evaluate \system\ on \subs{real-world datasets}{six datasets} with diverse textual and data characteristics, covering constraints ranging from simple checks to complex cross-record rules. \subs{Across constraint classification, extraction, and violation detection, \system\ outperforms baselines from data profiling and Text-to-SQL,}{\system\ outperforms baselines from data profiling and Text-to-SQL in constraint classification and extraction on all six datasets,} recovering constraints they miss and identifying real regulatory violations \rem{with high precision }in publicly released data\add{, including two large mortgage datasets with up to one million records}.







%Data quality constraints (aka data rules) are critical for high-quality, consistent data, yet datasets from data lakes, open data portals, and heterogeneous systems often lack them. Data profiling discovers constraints directly from data but cannot distinguish coincidental patterns from semantically meaningful rules. We pursue a different approach: extracting constraints for a dataset from the regulatory documents relevant to it. However, automated extraction is challenging because constraint sentences are mixed among definitions, procedures, and context, and regulatory text expresses data quality constraints with varied phrasing, implicit scope qualifiers, and conditional logic that must be correctly interpreted to produce executable constraints.
%
%
%
%We introduce the \emph{Text-to-Constraints} problem and present \system\ (\emph{\underline{Norma}tive Constraint Mapper}), a framework for extracting data quality constraints from regulatory text to validate datasets against their regulations. 
%To achieve high accuracy while scaling to large documents, \system\ decomposes the problem into constraint classification and constraint mapping, enabling implementations with different cost-accuracy trade-offs at each stage.
%The most effective configurations combine language models with novel techniques that link passages to schema-grounded constraints.
%We release annotated benchmarks that range from simple checks to complex business rules with conditionals and cross-record dependencies, along with metrics for extraction correctness and violation detection. Experiments show that \system\ outperforms baselines from data profiling, text-to-SQL, and data cleaning, extracting more accurate constraints and identifying more inconsistencies while scaling to large documents at low cost.



%Data quality constraints (or data rules) are key for reliable data, yet datasets from data lakes, open data sources, and heterogeneous systems often lack them.
%%Since manual definition is time-consuming and error-prone, 
%Data profiling aims to discover data rules directly from data; however, current solutions cannot distinguish coincidental patterns from meaningful data rules. We explore a different direction: many datasets have accompanying textual content, for instance, regulatory documents, policy manuals, and legal texts, which implicitly encode data rules. The challenge is that this kind of text does not explicitly state these rules and often mixes them with other content, making them easy to overlook.

%Data quality constraints (or rules) are essential for reliable data, yet datasets from data lakes, open data portals, and heterogeneous systems often lack them. Data profiling discovers constraints directly from data, but cannot distinguish coincidental patterns from semantically meaningful rules. We pursue a different direction: extracting constraints from text. Regulatory documents, policy manuals, and legal texts often define critical data quality rules but the challenge is how to extract,these they rarely explicitly state constraints, mixing them with other content and making extraction difficult.






%In this paper, we introduce the \emph{Text-to-Constraints} problem and present \system\ (\emph{\underline{Norma}tive Constraint Mapper}), a framework that, given a regulatory document and a related dataset, extracts data quality constraints that validate the data. We release annotated datasets targeting expressive constraints, ranging from simple checks to complex business rules with conditionals and multi-record dependencies. We also propose evaluation metrics for the correctness of constraint extraction and the effectiveness of constraint violation detection.
%To handle large documents and datasets, \system\ decomposes the problem into constraint classification and constraint mapping, and supports implementations with different cost-effectiveness trade-offs at each stage. 
%The most effective implementations leverage language models combined with constraint patterns grounded in the schema and data for better constraint extraction. \system\ outperforms diverse baselines in data profiling, text-to-SQL, and data cleaning, consistently extracting more accurate constraints and identifying more data inconsistencies, while scaling to larger documents and datasets at low cost.








